\chapter{Ratio, Root, and Strategy Tests}

\section{Ratio and Root Tests for Convergence}
\subsection{Absolute Convergence}
\index{absolute convergence}Suppose there is a series $\sum_{n=1}^\infty a_n$, 
then there is a corresponding series $\sum_{n=1}^\infty |a_n| = |a_1| + |a_2| 
+ |a_3| + \cdots$. If $\sum_{n=1}^\infty |a_n|$ is convergent, then the series 
$\sum_{n=1}^\infty a_n$ is called \textbf{absolutely convergent}. 

\textbf{Example}: Consider the alternating series 
$$\sum_{n=1}^\infty \frac{(-1)^{n-1}}{n^2} = 1 - \frac{1}{2^2} + 
\frac{1}{3^2} + \cdots$$\\
Is this series absolutely convergent?

\textbf{Solution}: We examine the corresponding series where we take the 
absolute value of each term:
$$\sum_{n=1}^\infty \left|\frac{(-1)^{n-1}}{n^2} \right| = \sum_{n=1}^
\infty \frac{1}{n^2}$$
We can identify $\sum_{n=1}^\infty \frac{1}{n^2}$ as a convergent $p$-series. 
Since $\sum_{n=1}^\infty \frac{1}{n^2}$ is convergent, we can state that 
$\sum_{n=1}^\infty \frac{(-1)^{n-1}}{n^2}$ is absolutely convergent. 

\textbf{Example} Is the convergent series $\sum_{n=1}^\infty 
\frac{(-1)^{n-1}}{n}$ absolutely convergent?

\textbf{Solution} We consider the sum of the absolute values of the terms:
$$\sum_{n=1}^\infty \left|\frac{(-1)^{n-1}}{n}\right| = \sum_{n=1}^\infty 
\frac{1}{n}$$\\
You should recognize this as the harmonic series, which is divergent. When a 
series is convergent but the corresponding series of absolute values is not, 
we call it \textbf{conditionally convergent}. \index{conditionally convergent}

We will not prove the theorem here, but it is useful to know that if a series 
$\sum_{n=1}^\infty a_n$ is absolutely convergent, then it is convergent. You 
can prove this yourself using the Comparison Test. 

\begin{Exercise}[label = absconv1]
Is the series given by 
$$\frac{\cos{1}}{1^2} + \frac{\cos{2}}{2^2} + \frac{\cos{3}}{3^3} + \cdots$$ \\
convergent or divergent?
\vspace{40mm}
\end{Exercise}

\begin{Answer}[ref = absconv1]
We can write the series as $\sum_{n=1}^\infty \frac{\cos{n}}{n^2}$. Since 
$n$ is real, we know that $n^2 > 0$ and we can say that $\sum_{n=1}^\infty 
\left| \frac{\cos{n}}{n^2} \right| = \sum_{n=1}^\infty \frac{|\cos{n}|}
{n^2}$. Additionally, $|\cos{n}| \leq 1$ for all $n$, and therefore 
$\frac{|\cos{n}|}{n^2} \leq \frac{1}{n^2}$. We know the series $\sum_{n=1}^
\infty \frac{1}{n^2}$ is convergent. And since we have shown that $\sum_{n=1}
^\infty \frac{|\cos{n}|}{n^2} \leq \sum_{n=1}^\infty \frac{1}{n^2}$, by the 
comparison test $\sum_{n=1}^\infty \frac{|\cos{n}|}{n^2}$ is convergent. 
Therefore, $\sum_{n=1}^\infty \frac{\cos{n}}{n^2}$ is absolutely convergent 
and therefore convergent. 
\end{Answer}

\begin{Exercise}[label = absconv2]
Determine whether each of the following series is absolutely or conditionally 
convergent. 
\begin{enumerate}
\item $\sum_{n=1}^\infty \frac{(-1)^n}{3n + 2}$
\item $\sum_{n=1}^\infty \frac{\sin{n}}{4^n}$
\item $\sum_{n=1}^\infty (-1)^{n-1} \frac{2n}{n^2 + 4}$ 
\end{enumerate}
\vspace{50mm}
\end{Exercise}

\begin{Answer}[ref = absconv2]
\begin{enumerate}
\item Conditionally Convergent. $\sum_{n=1}^\infty \left| \frac{(-1)^n}{3n+2} 
\right| = \sum_{n=1}^\infty \frac{1}{3n + 2}$ Applying the integral test to 
this sum: $\int_1^\infty \frac{1}{3x+2}\,dx = \lim_{t \to \infty} \int_1^t 
\frac{1}{3x+2}\,dx = \left[ \frac{1}{3} \ln{3x + 2} \right]_{x=1}^t = \lim_{t 
\to \infty} \left[ \ln{3x + 2} \right] - \ln{3(1) - 2} = \infty - 0 = \infty$. 
Since $\int_1^\infty \frac{1}{3x + 2}\,dx$ is divergent, $\sum_{n=1}^\infty 
\frac{1}{3n + 2}$ is divergent, and $\sum_{n=1}^\infty \frac{(-1)^n}{3n + 2}$ 
is conditionally convergent.
\item Absolutely Convergent. $\sum_{n=1}^\infty \left| \frac{\sin{n}}{4^n} 
\right| \leq \sum_{n=1}^\infty \frac{1}{4^n}$. Applying the integral test to 
$\sum_{n=1}^\infty \frac{1}{4^n}$: $\int_1^\infty \frac{1}{4^x}\,dx = \lim_{t 
\to \infty} \frac{-1}{4^x \ln{4}}|_{x=1}^t = \lim_{t \to \infty} \left[ \frac{
-1}{4^t \ln{4}} \right] - \frac{-1}{4^1 \ln{4}} = 0 + \frac{1}{4 \ln{4}} = 
\frac{1}{4 \ln{4}}$. Since $\int_1^\infty \frac{1}{4^x}\,dx$ is convergent, 
the series $\sum_{n=1}^\infty \frac{1}{4^n}$ is also convergent. And since 
$\sum_{n=1}^\infty \left| \frac{\sin{n}}{4^n} \right| \leq \sum_{n=1}^\infty 
\frac{1}{4^n}$, $\sum_{n=1}^\infty \left| \frac{\sin{n}}{4^n} \right|$ is 
also convergent, which shows that $\sum_{n=1}^\infty \frac{\sin{n}}{4^n}$ is 
absolutely convergent. 
\item Conditionally Convergent. We are asking if the series $\sum_{n=1}^\infty 
\left| (-1)^{n-1} \frac{2n}{n^2 + 4} \right|$ is convergent. $\sum_{n=1}^
\infty \left| (-1)^{n-1} \frac{2n}{n^2 + 4} \right| = \sum_{n=1}^\infty 
\frac{2n}{n^2 + 4}$ We will apply the Limit Comparison test and compare this 
series to the known, divergent series $\sum_{n=1}^\infty \frac{1}{n}$. $\lim_{
n \to \infty} \frac{\frac{2n}{n^2 + 4}}{\frac{1}{n}} = \lim_{n \to \infty} 
\frac{2n^2}{n^2 + 4} = 2 >0$. Therefore, by the Limit Comparison test, $\sum_{
n=1}^\infty \left| (-1)^{n-1} \frac{2n}{n^2 + 4} \right|$ is divergent AND 
$\sum_{n=1}^\infty (-1)^{n-1} \frac{2n}{n^2 + 4}$  is conditionally convergent. 
\end{enumerate}
\end{Answer}

\subsection{The Ratio Test}
\index{ratio test}The ratio test compares the $(n + 1)^{th}$ term of a series 
to the $n^{th}$ term and takes the limit as $n \to \infty$ of the absolute 
value of this ratio:
$$\lim_{n \to \infty} \left| \frac{a_{n + 1}}{a_n} \right| = L$$\\
There are three possible outcomes of the ratio test:
\begin{enumerate}
\item If $L < 1$, then the series $\sum_{n=1}^\infty a_n$ is absolutely 
convergent (and therefore convergent).
\item If $L = 1$, then the ratio test is inconclusive and we cannot draw any 
conclusions about whether $\sum_{n=1}^\infty a_n$ is convergent or divergent.
\item If $L > 1$ or $\lim_{n \to \infty} \frac{a_{n + 1}}{a_n} = \infty$, then 
the series $\sum_{n=1}^\infty a_n$ is divergent.
\end{enumerate}

\textbf{Example}: Apply the ratio test to determine if $\sum_{n=1}^\infty 
(-1)^n \frac{n^3}{3^n}$ is convergent or divergent. 

\textbf{Solution}: 
$$\lim_{n \to \infty} \left| \frac{\frac{(-1)^{n + 1}(n + 1)^3}{3^{n + 1}}}{
\frac{(-1)^n n^3}{3^n}} \right| = \lim_{n \to \infty} \frac{(n + 1)^3}{3^{n + 
1}} \cdot \frac{3^n}{n^3}$$
$$= \lim_{n \to \infty} \frac{(n + 1)^3}{n^3} \cdot \frac{3^n}{3 \cdot 3^n}$$
$$= \lim_{n \to \infty} \left( \frac{n + 1}{n} \right) \cdot \frac{1}{3} = 
\frac{1}{3} \lim_{n \to \infty} \left( \frac{n + 1}{n} \right) = \frac{1}{3} 
\cdot 1 = \frac{1}{3} $$
Since $L < 1$, the series $\sum_{n=1}^\infty (-1)^n \frac{n^3}{3^n}$ is 
absolutely convergent. 

The ratio test is most useful for series that contain factorials, constants 
raised to the $n^{th}$ power, or other products. 

\begin{Exercise}[label = ratio1]
[This question was originally presented as a multiple-choice, no-calculator 
problem on the 2012 AP Calculus BC exam.] Which of the following series are 
convergent?
\begin{enumerate}
\item $\sum_{n=1}^\infty \frac{8^n}{n!}$
\item $\sum_{n=1}^\infty \frac{n!}{n^{100}}$
\item $\sum_{n=1}^\infty \frac{n + 1}{(n)(n+2)(n+3)}$
\end{enumerate}
\end{Exercise}

\begin{Answer}[ref = ratio1]
Series 1 and 3 converge
\begin{enumerate}
\item We apply the ratio test: $\lim_{n \to \infty} \left| \frac{\frac{8^
{n + 1}}{(n + 1)!}}{\frac{8^n}{n!}} \right| = \lim_{n \to \infty} \frac{8 
\cdot 8^n}{(n + 1)(n!)} \cdot \frac{n!}{8^n} = \lim_{n \to \infty} 
\frac{8}{n + 1} = 0$. Therefore, the series converges. 
\item We apply the ratio test: $\lim_{n \to \infty} \left| \frac{\frac{(n + 1)!}
{(n + 1)^{100}}}{\frac{n!}{n^{100}}} \right| = \lim_{n \to \infty} \frac{(n + 1)
n!}{(n + 1)^{100}} \cdot \frac{n^{100}}{n!} = \lim_{n \to \infty} \left( 
\frac{n}{n + 1} \right)^{100} \cdot (n + 1) = \lim_{n \to \infty} \frac{n^{100}
}{(n + 1)^{99}} = \infty$. Therefore, the series diverges. 
\item We apply the comparison test: $\frac{n + 1}{(n)(n + 2)(n + 3)} = \frac{n}
{(n)(n + 2)(n + 3)} + \frac{1}{(n)(n + 2)(n + 3)} = \frac{1}{(n + 2)(n + 3)} + 
\frac{1}{(n)(n + 2)(n + 3)} = \frac{1}{n^2 + 5n + 6} + \frac{1}{n^3 + 5n^2 + 
6n} \leq \frac{1}{n^2} + \frac{1}{n^3}$. The series $\sum_{n = 1}^\infty 
\frac{1}{n^2}$ and $\sum_{n = 1}^\infty \frac{1}{n^3}$ are both convergent, 
because they are $p$-series with $p > 1$. Having established that $\sum_{n=1}
^\infty \frac{n + 1}{(n)(n + 2)(n + 3)} \leq \sum_{n=1}^\infty \frac{1}{n^2} 
+ \frac{1}{n^3}$ and that $\sum_{n=1}^\infty \frac{1}{n^2} + \frac{1}{n^3}$ 
converges, by the comparison test we can state that $\sum_{n=1}^\infty 
\frac{n + 1}{(n)(n+2)(n+3)}$ converges. 
\end{enumerate}
\end{Answer}

\begin{Exercise}[label = series2]
[This question was originally presented as a multiple-choice, calculator-allowed 
problem on the 2012 AP Calculus BC exam.]If the series $\Sigma_{n = 1}^\infty 
a_n$ converges and $a_n > 0$ for all $n$, which of the following statements 
must be true? Explain why. 
\begin{enumerate}
\item $\lim_{n \to \infty} \left| \frac{a_{n + 1}}{a_n} \right| = 0$
\item $|a_n|<1$ for all $n$
\item $\sum_{n=1}^\infty a_n = 0$
\item $\sum_{n=1}^\infty na_n$ diverges
\item $\sum_{n=1}^\infty \frac{a_n}{n}$ converges
\end{enumerate}
\end{Exercise}

\begin{Answer}[ref=series2]
\begin{enumerate}
\item This is not necessarily true. For a convergent series, the result of the 
ratio test is $L < 1$, so the limit could be $\neq 0$. 
\item This is not necessarily true. Consider the geometric series $\sum_{n=1}
^\infty 2(\frac{1}{2})^{n-1}$. This series is convergent because the common 
ratio is less than one, but the first term is $2(\frac{1}{2})^0 = 2 > 1$. 
\item This is not necessarily true. Again, consider the geometric series 
$\sum_{n=1}^\infty 2(\frac{1}{2})^{n-1}$, which converges to $4 \neq 0$. 
\item This is not necessarily true. Consider the $p$-series $\sum_{n=1}^
\infty \frac{1}{n^4}$. Then the series $\sum_{n=1}^\infty n \frac{1}{n^4} = 
\sum_{n=1}^\infty \frac{1}{n^3}$ is convergent. 
\item This must be true. By the comparison test, $\sum_{n=1}^\infty 
\frac{a_n}{n} \leq \sum_{n = 1}^\infty a_n$. Since $\sum_{n = 1}^\infty 
a_n$ converges, so much $\sum_{n = 1}^\infty na_n$. 
\end{enumerate}
\end{Answer}

\subsection{Root Test}
\index{root test}The root test examines the behavior of the $n^{th}$ root of 
$a_n$ as $n \to \infty$. Similar to the ratio test, there are three possible 
outcomes:
\begin{enumerate}
\item If $\lim_{n \to \infty} \sqrt[n]{|a_n|} = L < 1$, then the series $\sum
_{n=1}^\infty a_n$ is absolutely convergent, and therefore convergent.
\item If $\lim_{n \to \infty} \sqrt[n]{|a_n|} = L >1$ or $\lim_{n \to \infty} 
\sqrt[n]{|a_n|} = \infty$, then the series $\sum_{n=1}^\infty a_n$ is 
divergent.
\item If $\lim_{n \to \infty} \sqrt[n]{|a_n|} = L = 1$, then the Root Test is 
inconclusive. 
\end{enumerate}

The root test is best when there is a term or terms raised to the $n^{th}$ 
power. Consider the series $\sum_{n=1}^\infty \left( \frac{2n + 3}{3n + 2} 
\right)^n$:

\textbf{Example}: Is the series $\sum_{n=1}^\infty \left( \frac{2n + 3}{3n + 
2} \right)^n$ convergent or divergent?

\textbf{Solution}: Since $a_n$ consists of terms raised to the $n^{th}$ power, 
we will apply the root test for convergence:
$$\lim_{n \to \infty} \sqrt[n]{\left| \left( \frac{2n+3}{3n+2} \right)^n 
\right|} = \lim_{n \to \infty} \frac{2n+3}{3n+2} = \frac{2}{3} < 1$$ 
Therefore, by the root test, the series $\sum_{n=1}^\infty \left( \frac{2n + 
3}{3n + 2} \right)^n$ is convergent. 

\begin{Exercise}[label=root1]
Use the Root Test to determine whether the following series are convergent or 
divergent.
\begin{enumerate}
\item $\sum_{n=1}^\infty \left( \frac{3n^2 + 1}{n^2 - 4} \right)^n$
\item $\sum_{n=1}^\infty \frac{(-1)^{n}}{(\ln{n})^n}$
\item $\sum_{n=1}^\infty \left( 1+ \frac{1}{n} \right)^{n^2}$
\vspace{50mm}
\end{enumerate}
\end{Exercise}

\begin{Answer}[ref=root1] \begin{enumerate} \item $\lim_{n=1}^\infty \sqrt[n]{\left| \left( \frac{3n^2 + 1}{n^2 - 4} \right)^n 
\right|} = \lim_{n=1}^\infty \frac{3n^2 + 1}{n^2 - 4} = 3 >1$. Therefore, 
the series $\sum_{n=1}^\infty \left( \frac{3n^2 + 1}{n^2 - 4} \right)^n$ is 
divergent.
\item $\lim_{n \to \infty} \sqrt[n]{\left| \frac{(-1)^{n}}{(\ln{n})^n} \right|} 
= \lim_{n \to \infty} \sqrt[n]{\frac{1}{(\ln{n})^n}} = \lim_{n \to \infty} 
\frac{1}{\ln{n}} = \frac{1}{\infty} = 0 < 1$. Therefore, the series $\sum_
{n=1}^\infty \frac{(-1)^{n}}{(\ln{n})^n}$ is convergent.
\item $\lim_{n \to \infty} \sqrt[n]{\left| \left( 1 + \frac{1}{n} \right)^{n^2} 
\right|} = \lim_{n \to \infty} \left(1 + \frac{1}{n} \right)^n = e > 1$. 
Therefore, $\sum_{n=1}^\infty \left( 1+ \frac{1}{n} \right)^{n^2}$ is divergent. 
\end{enumerate}
\end{Answer}

\section{Strategies for Testing Series}\index{series!testing strategies}
When testing series for convergence, we want to choose a test based on the 
form of the series. While you may by tempted to try each test one-by-one until 
you find an answer, this quickly becomes cumbersome and time-consuming. 
Additionally, if you plan to take an AP Calculus exam, you need to be able to 
quickly choose an appropriate test as to conserve the time you have available 
for the exam. Here are some tips:
\begin{enumerate}
\item Check if the series is a $p$-series ($\sum_{n=1}^\infty \frac{1}{n^p}$). 
If so, then if $p > 1$, the series converges. Otherwise, the series diverges.
\item If the series is not a $p$-series, check to see if you can write it as a 
geometric series ($\sum_{n=1}^\infty ar^{n-1}$ or $\sum_{n=1}^\infty ar^n$). 
Recall that geometric series are convergent if $|r| < 1$ and divergent otherwise.
\item If the series cannot be written as a $p$-series or geometric series, but 
has a similar form, consider the comparison tests (the Direct Comparison Test 
and the Limit Comparison Test). When choosing a $p$-series to compare your 
series to, follow the guidelines outlined in the Comparison Tests section 
above. 
\item If you can see at a glance that $\lim_{n \to \infty} a_n \neq 0$, then 
apply the Test for Divergence to show the series is divergent. REMEMBER: $\lim_
{n \to \infty} a_n \neq 0$ implies the series $\sum_{n=1}^\infty a_n$ is 
divergent, but $\lim_{n \to \infty} a_n = 0$ \textbf{does not} necessarily imply the 
series $\sum_{n=1}^\infty a_n$ is convergent. 
\item If the series is alternating (has $(-1)^n$ or $(-1)^{n-1}$ in the term), 
the Alternating Series test may provide an answer. 
\item The Ratio Test is excellent for series with factorials, other products, 
or constants to the $n^{th}$ power. Remember that the Ratio Test will be 
inconclusive for $p$-series, rational functions of $n$, and algebraic functions 
of $n$. 
\item If $a_n$ is of the form $(b_n)^n$, use the Root Test. 
\item If $a_n = f(n)$ where $f(n)$ is continuous, positive, and decreasing and 
you can evaluate $\int_1^\infty f(x)\,dx$, use the Integral Test. 
\end{enumerate}

You do not need to treat this as a checklist, where you check for every 
condition. Rather, you should use this as a guide to quickly determine the 
convergence test most likely to be useful. 

\begin{Exercise}[label = test]
Choose an appropriate test to determine if the series is convergent of 
divergent. Apply the test and classify the series as convergent or divergent.
\begin{enumerate}
\item $\sum_{n=1}^\infty \frac{e^n}{n^2}$
\item $\sum_{n=1}^\infty \frac{3^n n^2}{n!}$
\item $\sum_{n=2}^\infty \frac{1}{n\sqrt{\ln{n}}}$
\item $\sum_{n=1}^\infty \left( \frac{n}{n+1} \right)^{n^2}$
\item $\sum_{n=1}^\infty \left( \sqrt[n]{2} - 1\right)^n$
\vspace{75mm}
\end{enumerate}
\end{Exercise}

\begin{Answer}[ref = test]
\begin{enumerate}
\item Divergent. Since there is a constant to the $n^{th}$ power and an algebraic 
function of $n$, we will try the Ratio Test. $\lim_{n \to \infty} \left| 
\frac{a_{n+1}}{a_n} \right| = \lim_{n \to \infty} \frac{e^{n + 1}}{(n + 1)^2} 
\cdot \frac{n^2}{e^n} = \lim_{n \to \infty} \frac{e^n \cdot e}{e^n} \cdot 
\left( \frac{n}{n + 1} \right)^2 = \lim_{n \to \infty} e \cdot \left( 
\frac{n}{n + 1} \right)^2 = e \cdot 1^2 = e >1$. Therefore, $\sum_{n=1}^\infty 
\frac{e^n}{n^2}$ is divergent. 
\item Convergent. Since there is a factorial, we will try the Ratio Test. $\lim_{n \to 
\infty} \left| \frac{a_{n+1}}{a_n} \right| = \lim_{n \to \infty} 
\frac{3^{n + 1}(n + 1)^2}{(n + 1)!} \cdot \frac{n!}{3^n n^2} = \lim_{n \to 
\infty} \frac{3 \cdot 3^n}{3^n} \cdot \frac{n!}{(n+1)n!} \cdot \left( 
\frac{n + 1}{n} \right)^2 = \lim_{n \to \infty} \frac{3(n+1)^2}{(n+1)n^2} = 
\lim_{n \to \infty} \frac{3(n+1)}{n^2} = 0 < 1$. Therefore, the series 
$\sum_{n=1}^\infty \frac{3^n n^2}{n!}$ is convergent. 
\item Divergent. Since $\int_2^\infty \frac{1}{x\sqrt{\ln{x}}}\,dx$ can be integrated, we 
will apply the integral test. $\int_2^\infty \frac{1}{x\sqrt{\ln{x}}}\,dx = 
\lim_{t \to \infty} \int_2^t \frac{1}{x\sqrt{\ln{x}}}\,dx$. Setting $u = 
\ln{x}$, then $du = \frac{dx}{x}$ and $\frac{1}{x\sqrt{\ln{x}}}dx = 
\frac{1}{\sqrt{u}}du$. Then we can say that $\int_2^\infty \frac{1}{x\sqrt{\ln{
x}}}\,dx = \lim_{t \to \infty} \int_{x=2}^{x=t} \frac{1}{\sqrt{u}}\,du = 
\lim_{t \to \infty} \left( \frac{-1}{2} \right) \sqrt{u} |_{x=2}^{x=t} = 
\lim_{t \to \infty} \left( \frac{-1}{2} \right) \sqrt{\ln{x}}_{2}^{t} = \left( 
\frac{-1}{2} \right) \lim_{t \to \infty} \sqrt{\ln{t}} - \left( \frac{-1}{2} 
\right) \sqrt{\ln{2}} = \infty$. Since the integral diverges, so does the 
series. 
\item Convergent. Since this series has terms to the $n^{th}$ power, we will 
try the Root Test. $\lim_{n \to \infty} \sqrt[n]{\left| \left( \frac{n}{n+1} 
\right)^{n^2} \right| } = \lim_{n \to \infty} \sqrt[n]{\left( \frac{n}{n+1} 
\right)^{n+2}} = \lim_{n \to \infty} \left(\frac{n}{n+1} \right)^{\frac{n^2}{
n}} = \lim_{n \to \infty} \left( \frac{n}{n+1} \right)^n = \lim_{n \to \infty} 
\left( \frac{1}{1 + \frac{1}{n}} \right)^n = \frac{1}{lim_{n \to \infty} \left(
1 + \frac{1}{n} \right)^n} = \frac{1}{e} < 1$ Therefore, by the root test, the 
series is convergent.
\item Convergent. This series also has terms raised to the $n^{th}$ power, we 
will try the Root Test again. $\lim_{n \to \infty} \sqrt[n]{\left| \left(\sqrt[
n]{2} - 1\right)^n \right| } = \lim_{n \to \infty} \sqrt[n]{\left( \sqrt[n]{2} 
- 1 \right)^n} = \lim_{n \to \infty}\left(\sqrt[n]{2} - 1 \right)^{n/n} = \lim_
{n \to \infty} \left( \sqrt[n]{2} - 1 \right) = \lim_{n \to \infty} 2^{1/n} - 
1 = 1 - 1 = 0 < 1$. Therefore, the series converges. 
\end{enumerate}
\end{Answer}
